\chapter*{Babagan Open Logic Project}
\addcontentsline{toc}{chapter}{Babagan Open Logic Project}


\textit{Open Logic Text} yaiku buku ajar sumber kabuka kang disusun
bebarengan, ngenani metalogika formal lan metodhe formal, wiwit tataran
madya (yaiku, sawise kuliah pambuka logika formal). Sanajan ditujokake
marang pamaca kang dhasar pasinaone dudu matematika (mligine mahasiswa
filsafat lan ilmu komputer), andharane tetep tliti.

Sawetara topik kang saiki wis dirembug bisa uga durung jangkep
andharane, lan akeh perangan isih mbutuhake revisi gedhe. Ing tembe,
kita ngrancang ngrembakakake teks iki supaya ngrembug luwih akeh topik.
Kita uga ngrancang nambahake glosarium, dhaptar wacan luwih lanjut,
cathetan sejarah, gambar, andharan kang luwih cetha, perangan kang
njlentrehake gayutane asil karo filsafat, ilmu komputer, lan matematika,
sarta luwih akeh soal lan tuladha ing teks iki.
Yen nemokake kaluputan utawa duwe usul,
\href{https://github.com/OpenLogicProject/OpenLogic/wiki/Contributing}{tulung ngandhani tim proyek}.

Proyek iki lumaku kanthi semangat sumber kabuka. Teks iki ora mung bisa
diakses kanthi bebas, nanging kita uga nyedhiyakake sumber LaTeX kanthi
Lisènsi Creative Commons Attribution. Lisènsi iki menehi hak marang
sapa wae kanggo ngundhuh, nggunakake, ngowahi, nata maneh, ngowahi
format, lan ngedum maneh karya kita, anggere menehi krédhit kang patut.
Katrangan liyane bisa diwaca ing situs Open Logic Project:
\href{http://openlogicproject.org/}{openlogicproject.org}.
